% A research alternative, not a checker-size claim.
% Requires amsmath, amssymb, amsthm; theorem/lemma/proposition/corollary defined.

\section{An alternative with deterministic, matrix-sensitive closure}
\label{sec:mixed-closure}

This section gives a different logical presentation of the same arithmetic.
It removes the independent case/progress choice from introduction by changing
the meaning and elimination of an atomic closure.  It does \emph{not} claim
a smaller complete checker: the additional arithmetic in elimination must
be counted.  In particular, the existing witness-selected introduction
already gives a deterministic introduction by another method.  The result
here is an alternative bi-interpretation theorem, not an optimality theorem.

The comparison is with the native scalar core in the pinned
\texttt{ConciseIPRA} snapshot.  Its existing progressive-binder proposal uses
progressive semantics at every matrix and a separate Boolean-cases principle.
The construction below uses progressive semantics only at literal equations,
retains anchored semantics at compound matrices, and derives numerical cases
through the embedding proved below.  No literature-wide priority claim is made.

\subsection{Complete native presentation}

Terms and formulas are finite syntax:
\[
 t::=x\mid0\mid St\mid I(n,a;u.h),\qquad
 B::=(s=t)\mid Q_yB\mid D_{x,y}B.
\]
Only $u$ is bound in $h$; $n$ and $a$ are outside that binding.
The displayed binders avoid capture, and contexts are finite reusable
lists of formulas.  Put $1=S0$ and introduce two metalinguistic
abbreviations, neither of which adds a term constructor:
\[
 A(n,t):=I(n,t;u.Su),\qquad
 d_B(n):=
 \begin{cases}
 n,&B\text{ is a literal equation},\\
 0,&B\text{ is compound}.
 \end{cases}
 \tag{MC1}\label{eq:mc-data}
\]
``Literal equation'' means that the outermost constructor of the fixed
matrix is equality.  An equation beneath $Q$ or $D$ is compound at that
outer level.  Numerical substitution preserves this classification.
This is a syntax-sensitive presentation, equivalently a grammar with
distinct atomic and compound closure forms.  It is not a claim that
one ordinary logical connective changes meaning under replacement by
an equivalent formula.  The rules below permit numerical substitution,
not unrestricted substitution of formulas for their matrices.

The two computation axioms, available in every context, are
\[
 I(0,a;u.h)=a,\qquad
 I(Sn,a;u.h)=h[I(n,a;u.h)/u].
 \tag{MC2}\label{eq:mc-comp}
\]
The remaining primitive rules are exactly
\[
 \frac{}{\Gamma\vdash B}\ (B\in\Gamma),
 \qquad
 \frac{\Gamma\vdash s=t\qquad\Gamma\vdash F(s)}
      {\Gamma\vdash F(t)},
 \tag{MC3}\label{eq:mc-hyp-tr}
\]
\[
 \frac{\Gamma\vdash Q_yB(y)\qquad
       \Gamma\vdash B(d_B(n))}
      {\Gamma\vdash B(A(d_B(n),t))},
 \tag{MC4}\label{eq:mc-qe}
\]
\[
 \frac{\Gamma\vdash D_{x,y}B(x,y)\qquad
       \Gamma,Q_yB(b,y)\vdash F}
      {\Gamma\vdash F},
 \tag{MC5}\label{eq:mc-de}
\]
\[
 \frac{\Gamma,B(t,d_B(z))\vdash B(t,Sz)}
      {\Gamma\vdash D_{x,y}B(x,y)}.
 \tag{MC6}\label{eq:mc-di}
\]
In transport, $F$ is any native formula template.  In
\eqref{eq:mc-de}, $b$ is fresh for $\Gamma$, $F$, and the passive
parameters of $B$.  In \eqref{eq:mc-di}, $z$ is fresh for $\Gamma$,
the witness $t$, and the passive parameters of $B$.  The displayed local
assumptions are discharged.  Call this calculus $\mathcal C_{\mathrm m}$.
There are no additional equality, successor, logical, or induction axioms.
As in the inherited core, the designated contradiction is $0=1$.

Rule \eqref{eq:mc-qe} has two numerical certificate parameters, $n,t$.
For compound $B$, its source is $B(0)$ and its target reduces by a
computation equation to $B(t)$.  For an equation $E$, it transports
progress from $E(n)$ to $E(A(n,t))$.  The order of arguments matters:
$A(0,t)=t$ is a direct computation, whereas $A(t,0)=t$ is a derived
induction theorem.  Reversing those arguments would invalidate the
bootstrap proof given next.

\subsection{Induction and the logical padding operation}

Weakening, renaming, ordinary-assumption cut, and capture-free numerical
substitution are admissible by induction on proofs.  For substitution,
the test defining $d_B$ is unchanged, and the computation equations are
schemas over fixed finite terms.  At eigenvariable rules, rename the
eigenvariable away from the substituted terms or inserted derivations.

The usual elementary equality rules are derived from computation and
transport.  For example, if $r=I(0,t;u.0)$, computation proves $r=t$;
transporting that equation along itself gives $t=t$.  The standard
equality motives then give symmetry, transitivity, and numerical congruence.

\begin{lemma}[Closure introduction and equation induction]
\label{lem:mc-ei}
The following are admissible:
\[
 \frac{\Gamma,B(d_B(z))\vdash B(Sz)}{\Gamma\vdash Q_yB(y)},
 \qquad
 \frac{\Gamma\vdash E(0)\qquad\Gamma,E(z)\vdash E(Sz)}
      {\Gamma\vdash E(t)},
\]
where the second rule requires a literal equation $E$, and the usual
eigenvariable conditions apply.  Also, for every native matrix $B$,
\[
 Q_yB(y),B(0)\vdash B(t).
 \tag{MC7}\label{eq:mc-zero-specialization}
\]
\end{lemma}
\begin{proof}
For closure introduction, apply \eqref{eq:mc-di} with an unused witness
coordinate and witness $0$, then eliminate the resulting $D$ formula.
Its opened $Q$ assumption is independent of the fresh witness, so it
is also the desired conclusion.

For \eqref{eq:mc-zero-specialization}, take $n=0$ in
\eqref{eq:mc-qe}.  Both possible matrices have $d_B(0)=0$.
Computation gives $A(0,t)=t$, and transport gives the stated target.
For equation induction, closure introduction applied to the supplied
step gives $Q_yE(y)$; use the supplied base and
\eqref{eq:mc-zero-specialization}.  No ordinary logical connective
or numerical-cases principle is used here.
\end{proof}

\begin{lemma}[A raw progress step]
\label{lem:mc-step}
For a literal equation $E$,
\[
 Q_yE(y),E(n)\vdash E(Sn).
\]
Moreover, $Q_yB(t,y)\vdash D_{x,y}B(x,y)$ for every matrix $B$.
\end{lemma}
\begin{proof}
Equation induction proves $A(n,1)=Sn$: the base is a computation,
and the successor computation reduces the induction step to congruence
of successor.  Use \eqref{eq:mc-qe} with the specified source $n$ and
with second parameter $1$, then transport along this equation.

For the second assertion use \eqref{eq:mc-di} with witness $t$.
For an equation matrix its premise follows from the raw progress step;
for a compound matrix it follows from
\eqref{eq:mc-zero-specialization} at $Sz$.
\end{proof}

We need a carefully restricted congruence fact.  If two matrices have
the \emph{same root category} (both literal equations or both compound)
and their instances are uniformly interderivable, then their $Q$
closures and their $D$ closures are interderivable.  For $Q$, use
closure introduction: convert its local antecedent to the old one,
use Lemma~\ref{lem:mc-step} in the atomic case or
\eqref{eq:mc-zero-specialization} in the compound case, and convert
the successor back.  For $D$, first open its witness and then repeat
that construction with the same witness.  The reverse proof is identical.
This assertion does \emph{not} permit replacing an atomic matrix by
an equivalent compound matrix beneath an already present closure.
Their closure meanings need not agree.

Define the selector term $C(v,a,b)=I(v,a;u.b)$, with unused $u$,
and the gate on formula syntax by
\[
 G_v(s=t):=(C(v,0,s)=C(v,0,t)),\quad
 G_v(Q_yB):=Q_yG_vB,\quad
 G_v(D_{x,y}B):=D_{x,y}G_vB.
\]
Bound variables are first renamed away from $v$.  Computation,
closure introduction, and the same-category congruence just proved
give, by external induction on finite formula syntax,
\[
 \vdash G_0B,\qquad G_{Sv}B\dashv\vdash B.
 \tag{MC8}\label{eq:mc-gate}
\]
The first induction uses the theorem for the smaller matrix to supply
the successor premise, ignoring its local assumption.  The second
uses the two selector computation equations at the atomic leaf and
same-category congruence at binders.  A gate preserves every root category.

\begin{lemma}[Compound padding]
\label{lem:mc-padding}
Let $r,v$ be fresh for $B$ and put
\[
                 P(B):=D_{r,v}G_vB.
\]
Then $P(B)$ is compound and $P(B)\dashv\vdash B$, uniformly in
every retained context and under numerical substitution.
\end{lemma}
\begin{proof}
Under the retained assumption $B$, introduce $P(B)$ with witness $0$.
Whether \eqref{eq:mc-di} supplies $G_zB$ or $G_0B$ as its local
assumption, its successor target $G_{Sz}B$ follows from the retained
$B$ by \eqref{eq:mc-gate}.  Conversely, open $P(B)$, use the theorem
$G_0B$ as the base of its $Q$ closure, and specialize at $1$ by
\eqref{eq:mc-zero-specialization}.  The result $G_1B$ gives $B$.
Fresh bound coordinates make these constructions commute with
capture-free numerical substitution, up to harmless renaming.
\end{proof}

\subsection{An embedding of the existing native scalar core}

Let $\mathcal C$ be the existing Eq/$Q$/$D$ core: its $Q$ elimination
is \eqref{eq:mc-zero-specialization}, and it has separate case and
literal-equation progress introductions for $D$.  Define a translation
that leaves all terms unchanged and pads every translated matrix:
\[
 (s=t)^\dagger=(s=t),\qquad
 (Q_yB)^\dagger=Q_yP(B^\dagger),\qquad
 (D_{x,y}B)^\dagger=D_{x,y}P(B^\dagger).
 \tag{MC9}\label{eq:mc-embed}
\]
All padding coordinates are fresh.  Contexts are translated pointwise.
Every translated binder has a compound matrix, even when its original
matrix was an equation.

\begin{theorem}[Proof translation]
\label{thm:mc-embed}
If $\Gamma\vdash_{\mathcal C}B$, then
$\Gamma^\dagger\vdash_{\mathcal C_{\mathrm m}}B^\dagger$.
The transformation is effective on finite derivations.
\end{theorem}
\begin{proof}
Induct on the derivation.  Hypotheses, computation, and transport are
preserved; the last assertion uses commutation of the translation
with numerical substitution.  Old $Q$ elimination first pads its
base proof, specializes the new compound closure, and unpads its
result.  Old $D$ elimination opens exactly the translated old $Q$
assumption.

For old case introduction, the translated old premise starts from
$B^\dagger(t,0)$ and gives $B^\dagger(t,Sz)$.  The new introduction
has compound matrix $P(B^\dagger)$: unpad its local base assumption,
insert the translated premise by cut, and pad the conclusion.

For old progress introduction with equation $E$, the translated
premise is $\Gamma^\dagger,E(t,z)\vdash E(t,Sz)$.  Introduce the
translated conclusion $D_{x,y}P(E(x,y))$ by the new compound case
rule.  Under its base $P(E(t,0))$, unpad to $E(t,0)$.
Lemma~\ref{lem:mc-ei}, with a separately fresh induction parameter,
gives $E(t,Sz)$; pad this result to finish the premise.  This is the
only old rule needing induction in the translation.
\end{proof}

The inherited scalar core reconstructs ordinary intuitionistic
first-order logic, numerical coverage, separation, and contextual
equation induction.  Theorem~\ref{thm:mc-embed} therefore imports
their proofs and their logical introduction/elimination rules for
translated formulas.  Those rules remain usable with arbitrary
additional retained new formulas: their proof transformations only
inspect the active assumptions, and weakening and cut were proved
above.  This observation suffices for the composite proof below;
it does not assume unrestricted cross-category binder congruence.

\subsection{Ordinary interpretation and both composites}

Let $T_{\mathrm{eq}}$ be the ordinary scalar realization of the existing
proposal: intuitionistic first-order equality logic over the iterator
terms, with \eqref{eq:mc-comp}, numerical coverage, separation
$(0=1)\to\bot$, and induction for literal equations with side assumptions.
The inherited arithmetic bridge identifies $T_{\mathrm{eq}}$ with
full first-order iPRA.  Our interpretation $\rho$ of the new calculus
in $T_{\mathrm{eq}}$ leaves terms and equality unchanged and sets
\[
 (Q_yB)^\rho=
 \begin{cases}
 \forall y\bigl(B^\rho(y)\to B^\rho(Sy)\bigr),
       &B\text{ is a literal equation},\\
 B^\rho(0)\to\forall y\,B^\rho(y),&B\text{ is compound},
 \end{cases}
 \qquad
 (D_{x,y}B)^\rho=\exists x\,(Q_yB)^\rho.
 \tag{MC10}\label{eq:mc-reading}
\]
Thus the atomic reading is genuinely changed; it is not the old
anchored reading with a proof choice silently removed.

\begin{lemma}[Soundness of the new rules]
\label{lem:mc-sound}
The interpretation $\rho$ sends new derivations to
$T_{\mathrm{eq}}$ derivations.
\end{lemma}
\begin{proof}
Only introduction and elimination merit comment.  Atomic introduction
is directly a proof of the displayed progressive statement with a
specified existential witness.  Compound introduction proves its
anchored statement by numerical coverage: zero is its retained base,
and successors are its supplied premise.

For atomic elimination, equation induction first proves in
$T_{\mathrm{eq}}$
\[
 A(n,0)=n,\qquad A(n,Sr)=S A(n,r).
\]
The first is induction on $n$.  The second is induction on $n$ with
$r$ as a passive parameter; both steps reduce by computation and
successor congruence.  Under the progressive premise and $E(n)$,
equation induction on $r$ now proves $E(A(n,r))$.  Taking $r=t$
gives \eqref{eq:mc-qe}.  The inducted formula is still one literal
equation, not an arbitrary compound formula.  Compound elimination
is ordinary anchored specialization, since $A(0,t)=t$.
The other rules have their immediate ordinary meanings.
\end{proof}

Let $J$ be the existing native interpretation of $T_{\mathrm{eq}}$
in $\mathcal C$, including its ordinary logical macros.  Define the
reverse interpretation $K$ by $K(\varphi)=J(\varphi)^\dagger$.
Its proof translation follows from Theorem~\ref{thm:mc-embed}.

\begin{theorem}[Whole-domain bi-interpretation]
\label{thm:mc-bi}
The interpretations $\rho,K$ are a bi-interpretation between
$\mathcal C_{\mathrm m}$ and $T_{\mathrm{eq}}$.  Both numerical
domains are whole, equality is actual equality, and both numerical
composites have the identity graph $x=y$.  Consequently the new
calculus is bi-interpretable and equiconsistent with iPRA.
\end{theorem}
\begin{proof}
First, for every old native formula $F$,
$\rho(F^\dagger)$ is provably equivalent to the old ordinary reading
of $F$.  Prove this externally by induction on $F$.
Every translated matrix is compound; hence its new closure has the
old anchored reading.  Lemma~\ref{lem:mc-padding}, transported by
Lemma~\ref{lem:mc-sound}, removes the inserted padding.  The old
ordinary composite theorem therefore gives
$\rho(K\varphi)\leftrightarrow\varphi$ for every ordinary formula
$\varphi$.

For the new composite, prove externally on native formulas that
$B\dashv\vdash K(B^\rho)$.  Equality is unchanged.  Suppose the
assertion is established for a matrix $B$ and write
$\widehat B=K(B^\rho)$; this is a formula in the image of $K$.

If $B=E$ is atomic, a new $Q_yE$ supplies the raw step
$E(z)\to E(Sz)$ by Lemma~\ref{lem:mc-step}.  Convert its two
instances to their hats and use the imported ordinary implication
and universal introductions.  This gives
$K(\forall y(E^\rho(y)\to E^\rho(Sy)))$.
Conversely, assume that imported progressive formula.  Under $E(z)$,
convert to $\widehat E(z)$, eliminate the imported universal and
implication, convert the result back to $E(Sz)$, and use derived
closure introduction.  These are exactly the two directions of the
atomic $Q$ composite.

For compound $B$, start instead with an assumption
$\widehat B(0)$, convert to $B(0)$, specialize the given new $Q_yB$
at an arbitrary fresh parameter, and convert back to its hat.
Imported implication and universal introduction yield the anchored
ordinary formula.  Conversely, under the base $B(0)$, that imported
formula supplies $B(Sz)$ after the same conversions; compound
closure introduction finishes.  No congruence is applied across
the potentially different roots of $B$ and $\widehat B$.

Finally, for $D_{x,y}B$, open its fresh witness, apply the just-proved
$Q$ equivalence for the matrix $B$, and introduce the imported
existential.  Conversely, eliminate that existential at a fresh
witness, convert its matrix back to $Q_yB$, and use the existential
injection derived in Lemma~\ref{lem:mc-step}.  All ordinary logical
rules here act on image formulas; arbitrary additional new assumptions
are merely retained.  This completes the formula composite proof.

Both interpretations leave every scalar term and numerical variable
unchanged, so their numerical composites are literally the identity.
Composition with the inherited $T_{\mathrm{eq}}$--iPRA
bi-interpretation yields the assertion.  The contradiction $0=1$
is unchanged, giving equiconsistency as well.
\end{proof}

\paragraph{Assessment.}
The result removes the independent introduction mode and exposes a
different distribution of arithmetic and logic: literal-equation
closures are progressive, while compound closures are anchored.
The padding embedding proves that this does not restrict ordinary
intuitionistic arithmetic.  However, elimination now builds an iterator
term and supplies an additional numerical parameter.  Counting only the
single introduction schema would conceal that cost.  No measured
improvement, proof-assistant verification, or minimality claim is made
for this alternative.
